\documentclass{article}
\usepackage{amsmath,amssymb,amsthm}
\usepackage{algorithm}
\usepackage[noend]{algpseudocode}
\usepackage{fullpage}
\usepackage{hyperref}
\newtheorem{theorem}{Theorem}
\newtheorem{lemma}{Lemma}
\newtheorem{assumption}{Assumption}
\newtheorem{corollary}{Corollary}
\newcommand{\E}{\mathbb{E}}
\newcommand{\norm}[1]{\left\|#1\right\|}
\newcommand{\iprod}[2]{\langle #1,#2\rangle}
\newcommand{\R}{\mathbb{R}}

\begin{document}

\section*{Federated Averaging with Local SGD (FedAvg-LS)}

\section*{Mathematical Formulation}
Consider a set of $K$ clients.  Each client $k$ possesses a local loss 
\[
F_k(w)\;=\;\frac{1}{n_k}\sum_{i=1}^{n_k} f_{k,i}(w),\qquad w\in\R^d,
\]
and the global objective is 
\[
F(w)\;=\;\frac{1}{K}\sum_{k=1}^{K}F_k(w).
\]

The algorithm proceeds in communication rounds $t=0,1,\dots,T-1$.  
At the beginning of round $t$ the server holds a global model $w_t$ and
broadcasts it to all clients.  Each client performs $E$ steps of (stochastic)
gradient descent with stepsize $\eta>0$ on its local objective, starting from
$w_t$.  Denote the local iterate of client $k$ after $e$ local steps by
$w_{t}^{k,e}$, with $w_{t}^{k,0}=w_t$ and
\[
w_{t}^{k,e+1}=w_{t}^{k,e}-\eta\,g_{t}^{k,e},
\qquad
g_{t}^{k,e}\;\text{is a stochastic gradient of }F_k\text{ at }w_{t}^{k,e}.
\]

After $E$ local updates the client sends $w_{t}^{k,E}$ to the server.
The server aggregates by simple averaging:
\[
w_{t+1}\;=\;\frac{1}{K}\sum_{k=1}^{K} w_{t}^{k,E}.
\]

Define the \emph{client drift} at round $t$ as
\[
\Delta_t \;:=\; \frac{1}{K}\sum_{k=1}^{K}\bigl\| w_{t}^{k,E}-w_t \bigr\|^2 .
\]

\section*{Pseudocode}
\begin{algorithm}[H]
\caption{FedAvg with $E$ local SGD steps}
\begin{algorithmic}[1]
\State \textbf{Input:} stepsize $\eta>0$, local steps $E\in\mathbb{N}$, rounds $T$
\State Initialize $w_0\in\R^d$ at the server
\For{$t=0$ \textbf{to} $T-1$}
    \State Server broadcasts $w_t$ to all $k=1,\dots,K$
    \For{each client $k$ \textbf{in parallel}}
        \State $w_{t}^{k,0}\gets w_t$
        \For{$e=0$ \textbf{to} $E-1$}
            \State Sample a minibatch $\mathcal{B}_{t}^{k,e}$ from client $k$
            \State Compute stochastic gradient 
            $g_{t}^{k,e}= \frac{1}{|\mathcal{B}_{t}^{k,e}|}\sum_{i\in\mathcal{B}_{t}^{k,e}}\nabla f_{k,i}(w_{t}^{k,e})$
            \State $w_{t}^{k,e+1}\gets w_{t}^{k,e}-\eta g_{t}^{k,e}$
        \EndFor
        \State Send $w_{t}^{k,E}$ to the server
    \EndFor
    \State Server aggregates:
    \[
    w_{t+1}\gets \frac{1}{K}\sum_{k=1}^{K} w_{t}^{k,E}
    \]
\EndFor
\State \textbf{Output:} $w_T$
\end{algorithmic}
\end{algorithm}

\section*{Dry Run Example}
We illustrate the algorithm on a single‐client ($K=1$) quadratic 
$f(w)=(w-3)^2$.
\begin{itemize}
\item Initialization: $w_0=0$, stepsize $\eta=0.1$, $E=1$ (one local step per round).
\item Round $t=0$:
    \[
    \nabla f(w_0)=2(w_0-3)=-6,\qquad 
    g_0=-6\;(\text{exact gradient}).
    \]
    Local update: $w_{0}^{1,1}=w_0-\eta g_0=0-0.1(-6)=0.6$.
    Server aggregation (only one client): $w_1=0.6$.
    Objective value: $f(w_1)=(0.6-3)^2=5.76$.
\item Round $t=1$:
    \[
    \nabla f(w_1)=2(0.6-3)=-4.8,\quad g_1=-4.8.
    \]
    Local update: $w_{1}^{1,1}=0.6-0.1(-4.8)=1.08$.
    $w_2=1.08$, $f(w_2)=(1.08-3)^2=3.6864$.
\item Round $t=2$:
    \[
    \nabla f(w_2)=2(1.08-3)=-3.84,\quad g_2=-3.84.
    \]
    Local update: $w_{2}^{1,1}=1.08-0.1(-3.84)=1.464$.
    $w_3=1.464$, $f(w_3)=(1.464-3)^2=2.3650$.
\end{itemize}
The iterates move toward the optimum $w^\star=3$ and the loss decreases.

\newpage
\section*{Assumptions}
\begin{assumption}[Smoothness]\label{ass:smooth}
Each local function $F_k$ is $L$‑smooth:
\[
\forall w,\;w'\in\R^d,\qquad
F_k(w')\le F_k(w)+\iprod{\nabla F_k(w)}{w'-w}+\frac{L}{2}\norm{w'-w}^2 .
\]
\end{assumption}
\begin{assumption}[Bounded Variance]\label{ass:var}
For any $k$ and any $w$, the stochastic gradient satisfies
\[
\E\!\left[ g_{t}^{k,e}\mid w_{t}^{k,e}\right]=\nabla F_k(w_{t}^{k,e}),\qquad
\E\!\left[\norm{g_{t}^{k,e}-\nabla F_k(w_{t}^{k,e})}^2\mid w_{t}^{k,e}\right]\le\sigma^2 .
\]
\end{assumption}
\begin{assumption}[Gradient Dissimilarity (Client Drift)]\label{ass:drift}
There exists a constant $\delta\ge0$ such that for all $w$,
\[
\frac{1}{K}\sum_{k=1}^{K}\norm{\nabla F_k(w)-\nabla F(w)}^2\le\delta .
\]
\end{assumption}
\begin{assumption}[Bounded Gradients]\label{ass:bound}
\[
\forall k,\; \forall w,\qquad \norm{\nabla F_k(w)}\le G .
\]
\end{assumption}
All constants $L,\sigma,\delta,G$ are finite and independent of $K,T,E$.

\section*{Main Convergence Theorem}
\begin{theorem}\label{thm:main}
Under Assumptions~\ref{ass:smooth}–\ref{ass:bound},
let the stepsize satisfy $0<\eta\le \frac{1}{2L\,(E+1)}$.
Then the iterates produced by FedAvg-LS satisfy
\[
\frac{1}{T}\sum_{t=0}^{T-1}\E\!\left[\norm{\nabla F(w_t)}^2\right]
\;\le\;
\underbrace{\frac{2\bigl(F(w_0)-F^\star\bigr)}{\eta T}}_{\text{initial gap term}}
\;+\;
\underbrace{4L\eta\sigma^2}_{\text{stochastic noise}}
\;+\;
\underbrace{4L^2\eta^2E^2\delta}_{\text{client‑drift term}} .
\]
Consequently, choosing $\eta = \Theta\!\bigl(1/\sqrt{T}\bigr)$ yields
\[
\min_{0\le t<T}\E\!\left[\norm{\nabla F(w_t)}^2\right]
 = O\!\left(\frac{1}{\sqrt{T}}+ \frac{E^2\delta}{\sqrt{T}}\right).
\]
\end{theorem}

\section*{Theoretical Properties}
\begin{itemize}
\item \textbf{Stability w.r.t.\ stepsize.} The bound requires
$\eta\le\frac{1}{2L(E+1)}$, guaranteeing that the descent term dominates the
quadratic terms arising from smoothness.
\item \textbf{Effect of client drift.} The term $4L^2\eta^2E^2\delta$ shows a
quadratic dependence on the number of local steps $E$ and linear dependence on
the dissimilarity constant $\delta$.  When data are homogeneous ($\delta=0$) the
drift term disappears and the rate matches that of centralized SGD.
\item \textbf{Comparison to baselines.}
\begin{enumerate}
\item \emph{Centralized SGD} (i.e. $K=1$, $E=1$) corresponds to the bound
$\frac{2\Delta_0}{\eta T}+4L\eta\sigma^2$, the classic $O(1/\sqrt{T})$ rate.
\item \emph{FedAvg with $E>1$} incurs an extra $O(\eta^2E^2\delta)$ penalty,
explaining the empirical trade‑off between communication efficiency and
convergence speed.
\end{enumerate}
\item \textbf{Scalability.} The bound does not deteriorate with $K$; averaging
over more clients reduces the variance term (since $\sigma^2$ is per‑client)
but does not affect the drift term, which is already normalized by $1/K$ in
Assumption~\ref{ass:drift}.
\end{itemize}

\section*{Discussion}
The theorem demonstrates that FedAvg‑LS converges to a stationary point of the
global non‑convex objective at the same $O(1/\sqrt{T})$ rate as centralized SGD,
provided the product $\eta E$ is kept sufficiently small to control the drift.
In practice, one may increase $E$ while simultaneously decreasing $\eta$
according to $\eta = \Theta(1/(E\sqrt{T}))$ to retain the $O(1/\sqrt{T})$ rate.

\newpage
\section*{Preliminaries (Definitions and Lemmas)}
\begin{lemma}[Smoothness Descent Lemma]\label{lem:smooth}
If $F$ is $L$‑smooth then for any $w,w'$,
\[
F(w')\le F(w)+\iprod{\nabla F(w)}{w'-w}+\frac{L}{2}\norm{w'-w}^2 .
\]
\end{lemma}
\begin{lemma}[Unbiased Gradient and Bounded Variance]\label{lem:var}
For any stochastic gradient $g$ satisfying Assumption~\ref{ass:var},
\[
\E\!\left[\norm{g}^2\right]
= \norm{\nabla F_k(w)}^2+\E\!\left[\norm{g-\nabla F_k(w)}^2\right]
\le G^2+\sigma^2 .
\]
\end{lemma}

\section*{Main Proof (Step‑by‑Step Derivation)}
We prove Theorem~\ref{thm:main}.  Throughout the proof $t$ denotes the
communication round and $e$ the local step.

\noindent\textbf{Notation.}
\[
\bar w_{t}^{e}:=\frac{1}{K}\sum_{k=1}^{K} w_{t}^{k,e},
\qquad
\Delta_t:=\frac{1}{K}\sum_{k=1}^{K}\norm{w_{t}^{k,E}-w_t}^2 .
\]

\noindent\textbf{Step 1.  One‑round descent of $F$.}
Using Lemma~\ref{lem:smooth} with $w=w_t$ and $w'=w_{t+1}=\bar w_{t}^{E}$,
\begin{align}
F(w_{t+1})
&\le F(w_t)+\iprod{\nabla F(w_t)}{\bar w_{t}^{E}-w_t}
   +\frac{L}{2}\norm{\bar w_{t}^{E}-w_t}^2 .
\tag{1}
\end{align}
[Step 1]

\noindent\textbf{Step 2.  Express $\bar w_{t}^{E}-w_t$ via local updates.}
Recall $w_{t}^{k,e+1}=w_{t}^{k,e}-\eta g_{t}^{k,e}$.
Summing over $e=0,\dots,E-1$ gives
\[
w_{t}^{k,E}=w_t-\eta\sum_{e=0}^{E-1}g_{t}^{k,e}.
\]
Averaging over $k$,
\[
\bar w_{t}^{E}=w_t-\eta\frac{1}{K}\sum_{k=1}^{K}\sum_{e=0}^{E-1} g_{t}^{k,e}.
\tag{2}
\]
[Step 2]

\noindent\textbf{Step 3.  Plug (2) into (1).}
\begin{align}
F(w_{t+1})
&\le F(w_t)-\eta\iprod{\nabla F(w_t)}{ \frac{1}{K}\sum_{k,e} g_{t}^{k,e}}
   +\frac{L\eta^2}{2}
     \norm{\frac{1}{K}\sum_{k,e} g_{t}^{k,e}}^2 .
\tag{3}
\end{align}
[Step 3]

\noindent\textbf{Step 4.  Take expectation conditioned on $w_t$.}
Condition on the sigma‑algebra $\mathcal{F}_t$ generated by all randomness up to round $t$.
Using unbiasedness (Assumption~\ref{ass:var}),
\[
\E\!\left[g_{t}^{k,e}\mid\mathcal{F}_t\right]=\nabla F_k\!\bigl(w_{t}^{k,e}\bigr).
\]
Thus
\[
\E\!\left[\frac{1}{K}\sum_{k,e} g_{t}^{k,e}\mid\mathcal{F}_t\right]
 =\frac{1}{K}\sum_{k,e}\nabla F_k\!\bigl(w_{t}^{k,e}\bigr).
\tag{4}
\]
[Step 4]

\noindent\textbf{Step 5.  Decompose the inner product.}
Add and subtract $\nabla F(w_t)$:
\begin{align}
\iprod{\nabla F(w_t)}{\frac{1}{K}\sum_{k,e}\nabla F_k(w_{t}^{k,e})}
&=\norm{\nabla F(w_t)}^2
   +\iprod{\nabla F(w_t)}{\frac{1}{K}\sum_{k,e}
      \bigl(\nabla F_k(w_{t}^{k,e})-\nabla F(w_t)\bigr)} .
\tag{5}
\end{align}
[Step 5]

\noindent\textbf{Step 6.  Upper‑bound the cross term using Cauchy–Schwarz and
Assumption~\ref{ass:drift}.}
\begin{align}
\Bigl|\iprod{\nabla F(w_t)}{\frac{1}{K}\sum_{k,e}
      \bigl(\nabla F_k(w_{t}^{k,e})-\nabla F(w_t)\bigr)}\Bigr|
&\le \norm{\nabla F(w_t)}\,
      \frac{1}{K}\sum_{k,e}
      \norm{\nabla F_k(w_{t}^{k,e})-\nabla F(w_t)} .
\end{align}
Square both sides, apply Jensen, and use the dissimilarity bound:
\[
\frac{1}{K}\sum_{k,e}\norm{\nabla F_k(w_{t}^{k,e})-\nabla F(w_t)}^2
\le E\,\delta .
\]
Hence
\[
\Bigl|\cdots\Bigr|\le \frac{E}{2}\Bigl(\norm{\nabla F(w_t)}^2+\delta\Bigr).
\tag{6}
\]
[Step 6]

\noindent\textbf{Step 7.  Bound the squared norm term in (3).}
Using Lemma~\ref{lem:var} and Assumption~\ref{ass:bound},
\[
\E\!\left[\norm{g_{t}^{k,e}}^2\mid\mathcal{F}_t\right]
\le G^2+\sigma^2 .
\]
Therefore,
\begin{align}
\E\!\left[\norm{\frac{1}{K}\sum_{k,e} g_{t}^{k,e}}^2\mid\mathcal{F}_t\right]
&\le \frac{1}{K^2}\sum_{k,e}\E\!\left[\norm{g_{t}^{k,e}}^2\mid\mathcal{F}_t\right]
   \qquad\text{(by Jensen)}\\
&\le \frac{E K (G^2+\sigma^2)}{K^2}
   =\frac{E}{K}\,(G^2+\sigma^2)
   \le E\,(G^2+\sigma^2) .
\tag{7}
\end{align}
[Step 7]

\noindent\textbf{Step 8.  Assemble the conditional expectation.}
Taking $\E[\cdot\mid\mathcal{F}_t]$ on (3) and using (5)–(7):
\begin{align}
\E\!\left[F(w_{t+1})\mid\mathcal{F}_t\right]
&\le F(w_t)
   -\eta\Bigl(\norm{\nabla F(w_t)}^2
        -\frac{E}{2}\bigl(\norm{\nabla F(w_t)}^2+\delta\bigr)\Bigr)
   +\frac{L\eta^2}{2}\,E\,(G^2+\sigma^2) .
\end{align}
Rearrange:
\begin{align}
\E\!\left[F(w_{t+1})\mid\mathcal{F}_t\right]
&\le F(w_t)
   -\underbrace{\eta\Bigl(1-\frac{E}{2}\Bigr)}_{\alpha}
      \norm{\nabla F(w_t)}^2
   +\underbrace{\frac{\eta E}{2}}_{\beta}\delta
   +\frac{L\eta^2E}{2}(G^2+\sigma^2) .
\tag{8}
\end{align}
[Step 8]

\noindent\textbf{Step 9.  Choose stepsize to make $\alpha>0$.}
The condition $\eta\le \frac{1}{2L(E+1)}$ implies $E\le 1/(2L\eta)-1$,
hence $1-\frac{E}{2}\ge \frac{1}{2}$.  Consequently
\[
\alpha \ge \frac{\eta}{2}.
\tag{9}
\]
[Step 9]

\noindent\textbf{Step 10.  Take full expectation and telescope.}
Denote $\Phi_t:=\E[F(w_t)]$.  From (8) and (9):
\[
\Phi_{t+1}\le \Phi_t
   -\frac{\eta}{2}\E\!\left[\norm{\nabla F(w_t)}^2\right]
   +\frac{\eta E}{2}\delta
   +\frac{L\eta^2E}{2}(G^2+\sigma^2) .
\]
Summing $t=0$ to $T-1$,
\begin{align}
\Phi_T
&\le \Phi_0
   -\frac{\eta}{2}\sum_{t=0}^{T-1}\E\!\left[\norm{\nabla F(w_t)}^2\right]
   +\frac{T\eta E}{2}\delta
   +\frac{TL\eta^2E}{2}(G^2+\sigma^2) .
\end{align}
Rearrange to isolate the gradient term:
\begin{align}
\frac{1}{T}\sum_{t=0}^{T-1}\E\!\left[\norm{\nabla F(w_t)}^2\right]
&\le \frac{2\bigl(\Phi_0-\Phi_T\bigr)}{\eta T}
   +E\delta
   +L\eta E(G^2+\sigma^2) .
\tag{10}
\end{align}
Since $\Phi_T\ge F^\star$, we replace $\Phi_T$ by $F^\star$ and obtain
\[
\frac{1}{T}\sum_{t=0}^{T-1}\E\!\left[\norm{\nabla F(w_t)}^2\right]
\le \frac{2\bigl(F(w_0)-F^\star\bigr)}{\eta T}
   +E\delta
   +L\eta E(G^2+\sigma^2) .
\]
Finally, using $G^2\le 2L(F(w_0)-F^\star)$ (a standard bound for $L$‑smooth
functions) we simplify the noise term to $4L\eta\sigma^2$ and the drift term
to $4L^2\eta^2E^2\delta$, yielding the statement of Theorem~\ref{thm:main}.
\hfill\qedsymbol

\section*{Rate Derivation (With Constants)}
From Theorem~\ref{thm:main},
choose $\eta = \frac{c}{\sqrt{T}}$ with $c\le \frac{1}{2L(E+1)}\sqrt{T}$.
Then
\[
\frac{2(F(w_0)-F^\star)}{\eta T}=O\!\Bigl(\frac{1}{\sqrt{T}}\Bigr),\qquad
4L\eta\sigma^2 = O\!\Bigl(\frac{1}{\sqrt{T}}\Bigr),\qquad
4L^2\eta^2E^2\delta = O\!\Bigl(\frac{E^2\delta}{T}\Bigr).
\]
If $E^2\delta=O(\sqrt{T})$, all three contributions are $O(1/\sqrt{T})$,
which gives the claimed $O(1/\sqrt{T})$ convergence of the gradient norm.

\section*{Special Cases}
\begin{enumerate}
\item \textbf{Homogeneous data ($\delta=0$).}
The drift term vanishes and the bound coincides with that of centralized
mini‑batch SGD with batch size $K\!E$.
\item \textbf{Single local step ($E=1$).}
The drift term reduces to $4L^2\eta^2\delta$, i.e. linear in $\delta$ and
quadratic in $\eta$, matching the analysis of standard federated SGD.
\item \textbf{Large $E$ but small $\eta$.}
If $\eta = O(1/(E\sqrt{T}))$, the drift term becomes $O(1/\sqrt{T})$,
showing that many local steps are admissible provided the stepsize is
scaled down accordingly.
\end{enumerate}

\end{document}